\section{From one run to the spread of its result}
\label{sec:seeds}

Fix a policy $\theta_0$ and start two runs from it. They see the same prompts and the same
objective, and differ only in the randomness of their rollouts. This is the comparison a
practitioner makes when they change a seed and rerun, and the question this section answers is how
far apart the numbers they report end up --- using only one of the runs.

\subsection{Noise is filtered, not accumulated}

Write $\Delta_t := \theta^A_t - \theta^B_t$. With updates $\theta \leftarrow \theta + \eta\ghat$
and $J_t := \partial\gbar/\partial\theta$ the Jacobian of the mean update,
\begin{equation}
  \label{eq:seed-recursion}
  \Delta_{t+1} = A_t \Delta_t + \eta\,\nu_t,
  \qquad A_t := I + \eta J_t,
\end{equation}
to first order in $\Delta$, with $\Cov(\nu_t) = 2\Sigma_t/P$. The batch noise $\Sigma_t$ is left abstract here and made
explicit in Section~\ref{sec:theory}, where it is computed exactly at finite group size; nothing
in this section depends on which estimator produced it.

If the objective were flat, $J = 0$, this would be a random walk: divergence would accumulate,
$\E\|\Delta_t\|^2 \propto t$, and a longer run would be strictly less reproducible than a short one.
That is the implicit model behind treating run-to-run variance as something to be measured by
brute force. But two RLVR runs are climbing the same objective, and where it is locally concave
$A_t$ is a contraction: divergence injected at one update decays over the ones that follow.

The consequence is that noise from update $t$ is not added to the final divergence. It is passed
through every subsequent update first.

\begin{theorem}[Finite-time covariance]
  \label{thm:propagation}
  Let $S_t$ denote the covariance of a run about its mean trajectory, with $S_0 = 0$ when both runs
  start from the same policy. Under the linearisation \eqref{eq:seed-recursion},
  \begin{equation}
    \label{eq:lyapunov}
    S_{t+1} = A_t S_t A_t^{\!\top} + Q_t,
    \qquad Q_t := \eta^2 \Sigma_t / P,
  \end{equation}
  and therefore
  \begin{equation}
    \label{eq:unrolled}
    S_T = \sum_{t=0}^{T-1} \Phi_{T,t+1}\, Q_t\, \Phi_{T,t+1}^{\!\top},
    \qquad \Phi_{T,s} := \!\!\prod_{u=s}^{T-1}\!\! A_u ,
  \end{equation}
  the product taken in order of decreasing $u$.
  Two independent seeds satisfy $\Cov(\theta^A_T - \theta^B_T) = 2S_T$, so their expected policy
  divergence is
  \begin{equation}
    \label{eq:predicted-kl}
    \E\big[\KL(\pi^A \,\|\, \pi^B)\big] \;\approx\; \tr\!\big(F_T\, S_T\big).
  \end{equation}
\end{theorem}

\begin{proof}
  Taking the covariance of \eqref{eq:seed-recursion} for a single run about its mean and using
  independence of $\nu_t$ from $\Delta_t$ gives \eqref{eq:lyapunov}; unrolling it from $S_0 = 0$
  gives \eqref{eq:unrolled}. Two independent runs have independent noise, so the covariance of
  their difference is twice that of one about the mean. Expanding
  $\KL(\pi^A\|\pi^B) \approx \tfrac{1}{2}\Delta^{\!\top}F\Delta$ and taking expectations gives
  $\tfrac{1}{2}\tr(F\cdot 2S_T)$.
\end{proof}

Equation \eqref{eq:unrolled} is the object this section is about, and it is worth reading as a
convolution: what survives to the end is the noise injected at each update, weighted by how
strongly the remaining trajectory preserves a perturbation entering there. Write
\begin{equation}
  \label{eq:kernel}
  k_t \;:=\; \tr\!\big(F_T\, \Phi_{T,t+1} Q_t \Phi_{T,t+1}^{\!\top}\big),
  \qquad
  \E[\KL] = \sum_{t} k_t,
\end{equation}
for the contribution of update $t$. Two runs are not equally sensitive to noise at every point in
training, and $k_t$ says where the sensitivity is.

\subsection{Why a single timescale is not enough}

An earlier version of this work replaced the spectrum of $A_t$ by one effective eigenvalue,
giving $\E[\KL_\infty] = \tau_c D_{\mathrm{noise}}$ for a persistence time $\tau_c$. That
approximation is the natural first move and it fails badly:
Section~\ref{sec:exp-seeds} measures it at a median error of two orders of magnitude, against
$9\%$ for \eqref{eq:unrolled}. The reason is visible in \eqref{eq:unrolled}. $A_t$ has a spectrum,
fast directions equilibrate within a few updates while slow ones set the final level, and a single
fitted exponential is pinned by the early rise and blind to what dominates the answer. We report
the scalar law only as the comparator it turned out to be.

\subsection{What the recursion predicts}

\paragraph{Divergence stops growing when, and because, the run is learning.} Contraction requires
$J_t$ to be meaningfully negative definite, which requires a signal. With no learning signal
\eqref{eq:seed-recursion} degenerates to a random walk. This is testable directly by removing the
signal while leaving the noise untouched, which Section~\ref{sec:exp-seeds} does.

\paragraph{The level is set by batch composition.} $Q_t \propto \Sigma_t/P$, so halving the
injected noise requires quadrupling the batch, and nothing in \eqref{eq:unrolled} depends on the
number of updates once the sum has converged.

\paragraph{A memory horizon, if contraction is strong enough to make one.} If $\Phi_{T,t+1}$
decays with $T - t$, the sum in \eqref{eq:unrolled} is dominated by its last terms and the final
$h \ll T$ updates account for nearly all of the spread; reproducing a result would then mean
reproducing its end. Where contraction is weak, $k_t$ is flat and every update counts the same.
Section~\ref{sec:exp-seeds} asks which of these holds and finds the second, over every setting we
can reach.

\subsection{Carrying it to the number that gets reported}
\label{sec:metric}

Equation \eqref{eq:predicted-kl} is a statement about policies. What a paper reports is a score:
pass rate, exact match, win rate, mean reward. Write $M(\theta)$ for that score as a function of
the parameters. Linearising it about the endpoint the run reached,
\begin{equation}
  \label{eq:metric-variance}
  \Var_{\text{seed}}\big[M(\theta_T)\big] \;\approx\; \nabla M(\theta_T)^{\!\top} S_T\, \nabla M(\theta_T),
\end{equation}
which needs $S_T$ --- a matrix nobody can form for a real model. It is not needed. Substituting
\eqref{eq:unrolled} turns the quadratic form into a sum of scalars.

\begin{theorem}[Adjoint form of the reported spread]
  \label{thm:adjoint}
  Define $b_T := \nabla M(\theta_T)$ and $b_s := A_s^{\!\top} b_{s+1}$ for $s = T-1, \dots, 0$.
  Then, under the assumptions of Theorem~\ref{thm:propagation},
  \begin{equation}
    \label{eq:adjoint}
    \begin{aligned}
      \Var_{\text{seed}}\big[M(\theta_T)\big]
      &\;=\; \sum_{t=0}^{T-1} b_{t+1}^{\!\top} Q_t\, b_{t+1} \\
      &\;=\; \frac{\eta^2}{P}\sum_{t=0}^{T-1} \Var\!\big(b_{t+1}^{\!\top}\hat{g}_t\big),
    \end{aligned}
  \end{equation}
  where the variance on the right is over one batch at $\theta_t$.
\end{theorem}

\begin{proof}
  $b_s = \Phi_{T,s}^{\!\top}\nabla M$ by induction. Substituting \eqref{eq:unrolled} into
  \eqref{eq:metric-variance} gives $\sum_t (\Phi_{T,t+1}^{\!\top}\nabla M)^{\!\top} Q_t
  (\Phi_{T,t+1}^{\!\top}\nabla M)$. The second equality is $Q_t = \eta^2\Sigma_t/P$ and
  $b^{\!\top}\Sigma b = \Var(b^{\!\top}\hat{g})$.
\end{proof}

Every term of \eqref{eq:adjoint} is the variance of a batch gradient projected onto a single
direction --- one scalar per update, which a trainer can accumulate alongside the update it is
already making. What the backward pass costs depends on $A_s^{\!\top}$, and here the estimator
family pays for itself.

\begin{proposition}[The mean update is a gradient field]
  \label{prop:potential}
  For any count-based estimator, $\gbar(\theta) = \nabla_\theta\,\Phi(\theta)$ with
  $\Phi(\theta) = \E_x[\Lambda(p_x)]$ and $\Lambda' = \lambda$. Hence $J = \nabla^2\Phi$ is
  symmetric and $A_s^{\!\top} = A_s$.
\end{proposition}

\begin{proof}
  By Proposition~\ref{prop:mean}, proved in Section~\ref{sec:theory}, the per-prompt mean is
  $\lambda(p_x, G)\,\nabla_\theta p_x$, and
  $\lambda$ depends on $\theta$ only through $p_x$. So
  $\lambda(p_x)\nabla p_x = \nabla \Lambda(p_x)$, and averaging over $x$ commutes with the
  gradient.
\end{proof}

The backward pass is therefore a Hessian-vector product, which finite differences deliver from two
gradient evaluations and reverse mode from one double backward. Neither $J$ nor $S_T$ is ever
formed. For RLOO, where $\lambda \equiv 1$, the potential is the pass rate itself: the mean update
ascends exactly the quantity being reported.

Two consequences worth stating plainly. First, this is a \emph{forecast}: \eqref{eq:adjoint} reads
only the states one run visited, so it can be computed before any second run exists, which is how
Section~\ref{sec:exp-seeds} tests it. Second, restricting $\Sigma_t$ to one of its two terms
attributes the answer by origin --- $\Sigma_b$ is which prompts were drawn, $\Sigma_w$ is which
responses were sampled --- and the two shares are exact and sum to one.

\subsection{Where the argument stops}

Everything above linearises the update around the shared trajectory, which requires the two runs to
stay close enough for $\gbar$ to be locally linear between them, and requires
$\eta\|J_t\| \ll 1$. Both fail when the step grows relative to the curvature, and RLVR supplies a
concrete way for that to happen: a controller holding drift at a fixed target must raise $\eta$ as
the signal decays, and $\E_x[p(1-p)] \to 0$ as a task is solved takes the signal with it.
Section~\ref{sec:exp-seeds} shows runs crossing out of the regime exactly this way.
